\documentclass[10pt,a4paper]{article}
\usepackage{amsmath,amssymb,amsthm,graphicx}
\title{Statistical-zero-knowledge containment from low-error randomized many-one reductions to random strings}
\author{Michael Saks, Rahul Santhanam}
\date{doi:10.4230/LIPIcs.CCC.2022.29; CC BY4.0}
\begin{document}
\maketitle
\noindent\textit{Editorial note.} Original authors: Michael Saks and Rahul Santhanam. Official CCC2022 article29 (2022-07-11), DOI10.4230/LIPIcs.CCC.2022.29. Creative Commons Attribution4.0.\par\medskip
\noindent\textit{Editorial note.} License: https://creativecommons.org/licenses/by/4.0/\par\medskip
\noindent\textit{Editorial note.} Faithful editable TeX transcription of formal author Theorem39 and its complete AppendixA proof pp29:29-29:30, exact universal-machine/random-string promise definitions and standard symmetry-of-information statement, complete published bibliography. Original author mathematical expressions retained; no corrections or additional proofs. This formal low-error many-one SZK containment is independent from the general nonadaptive AM containment; informal introduction statements are not used as targets.\par\medskip
\noindent\textit{Allowed standard prerequisites (editorial).}\begin{itemize}
\item Unbounded Kolmogorov symmetry of information. Exact author Proposition11 retained.
\item Entropy Difference in SZK and closure under polynomial many-one reductions. The complete standard containment fact as stated by the authors in AppendixA is retained in the proof.
\end{itemize}\medskip
\paragraph{Author descriptions and Kolmogorov complexity (Section 2.4).}
Throughout we assume a fixed pairing function $\langle\cdot\rangle$ and a standard encoding of Turing machines $M$ by binary strings, with $\widetilde M$ denoting the encoding of $M$.
We fix a universal Turing machine $U$ that for any Turing machine $M$ and input $x$ takes as input $\langle\widetilde M,x\rangle$ and outputs $M(x)$. The Kolmogorov complexity $K(q)$ of $q$ (with respect to $U$), denoted $K(q)$, is the minimum length of $y$ such that $U(y)=q$. We define the Kolmogorov complexity of a machine $M$ to be $K(\widetilde M)$.
\paragraph{Proposition 11 (Symmetry of information).}
Let $x,y$ be strings. Then $K(x;y)=K(x)+K(y\mid x)+O(\log(K(x,y)))$.
\paragraph{Author random-string promise problems (Section 2.5).}
In the $R_K$ problem, the instance is a string $x$, and the question is whether $K(x)\geq|x|/2$. Given a function $g:\mathbb N\to\mathbb N$, the $R_K$ problem with gap $g$ is the promise problem whose NO instances are strings $x$ with $K(x)\geq|x|/2$ and whose YES instances are strings $x$ with $K(x)<|x|/2-g(|x|)$. Given a function $\alpha:\mathbb N\to\mathbb R$ such that $\alpha(n)\geq1$ for all $n$, the $R_K$ problem with multiplicative gap $\alpha$ is the promise problem whose NO instances are strings $x$ with $K(x)\geq|x|/2$ and whose YES instances are strings $x$ such that $K(x)<|x|/(2\alpha(|x|))$.

\paragraph{Theorem 39.}
Let $L$ be any decidable language such that there is a randomized poly-time $m$-reduction with error $1/n^{\omega(1)}$ from $L$ to $R_K$ with gap $g$ for some time-constructible $g=\omega(\log(n))$. Then $L\in\mathsf{SZK}$.

\paragraph{Appendix A. Limitations of Randomized $m$-Reductions to Approximating Kolmogorov Complexity.}
In this section, we show that any language that has a low-error randomized polynomial-time $m$-reduction to a gap version of the Kolmogorov random strings is solvable in Statistical Zero Knowledge.
\begin{proof}
Let $L$ be a language as in the statement of the theorem. Let $f$ be a randomized poly-time $m$-reduction with error $1/n^{\omega(1)}$ to $R_K$ with gap $g$ for some $g=\omega(\log(n))$.
We can assume without loss of generality that any query made by $f$ on input $x$ of length $n$ is of length $p(n)$ for some polynomial $p$. Indeed, let $p(n)$ be an upper bound on the running time of $f$. We modify the reduction by ``padding'' any query made by $f$ to length $m=2p(n)$ as follows. Let $c$ be the constant in the $O(\cdot)$ in Proposition 11. If the query $q$ is of length $k$, we sample a string $y=r0^{(m-k)/2-2c\log(n)}$ where $r\in\{0,1\}^{(m-k)/2+2c\log(n)}$ is chosen uniformly at random, and output the query $qy$. Observe that by Proposition 11, with probability $1-1/n^{\omega(1)}$, if $q$ is a YES instance of $R_K$, then $qy$ is a YES instance of $R_K$, and if $q$ is a NO instance of $R_K$, then $qy$ is a NO instance of $R_K$.
For each $x\in\{0,1\}^*$, $f(x)$ is a random variable supported on $\{0,1\}^{2p(|x|)}$. We claim that the following holds. If $x\notin L$, the random variable $f(x)$ has entropy at least $p(|x|)-g(2p(|x|))/2+1$, and if $x\in L$, the random variable $f(x)$ has entropy at most $p(|x|)-g(2p(|x|))/2-1$.
We first show that the claim implies that $L\in\mathsf{SZK}$, and then establish the claim. To see that the claim implies $L\in\mathsf{SZK}$, we use the fact that Entropy Difference is in $\mathsf{SZK}$. Namely, there is a statistical zero-knowledge protocol that, given a circuit $C$ sampling a distribution over $m$ bits and a parameter $s$, accepts when the entropy of $C$ is at least $s+1$ and rejects when the entropy of $C$ is at most $s-1$. Note that a circuit $C$ sampling $f(x)$ can easily be computed from $x$. Since $\mathsf{SZK}$ is closed under poly-time $m$-reductions, we have that $L\in\mathsf{SZK}$, by reducing to Entropy Difference with parameter $m/2-g(m)/2$.
Next we establish the claim. Assume for the sake of contradiction that the claim fails for infinitely many inputs. Let $\{n_i\}$ be an infinite sequence of numbers such that the claim fails for at least one input of length $n_i$ for each $i$, and let $x_i$ be the first input of length $n_i$ for which the claim fails. Note that $K(x_i)=O(\log(n_i))$. The reason is that we can describe $x_i$ by a program which has $n_i$ in binary encoded into it, and then searches for the first string $z$ of length $n_i$ such that the claim fails for $z$, i.e., either $z\in L$ and the entropy of $f(z)$ is greater than $p(|z|)-g(2p(|z|)/2-1$, or $z\notin L$ and the entropy of $f(z)$ is smaller than $p(|z|)-g(2p(|z|))/2+1$.
Let $i$ be large enough and $x_i$ be the first string of length $n_i$ for which the claim fails. Let $m=2p(n_i)$. Either $x_i\in L$ and the entropy of $f(x_i)$ is greater than $m/2-g(m)/2-1$, or $x_i\notin L$ and the entropy of $f(x_i)$ is smaller than $m/2-g(m)/2+1$. We show that either case leads to a contradiction.
In the first case, the random variable $f(x_i)$ has entropy greater than $m/2-g(m)/2-1$. Then with probability at least $1/m$ over $y\in\{0,1\}^m$ sampled from $f(x_i)$, $K(y)>m/2-g(m)$. Indeed, if not, the entropy of $f(x_i)$ is at most
\[\frac1m*m+(1-1/m)*(m/2-g(m))\leq m/2-g(m)+1<m/2-g(m)/2-1,\]
which contradicts the lower bound on the entropy of $f(x_i)$. But if $K(y)>m/2-g(m)$ with probability at least $1/m$, the reduction cannot be correct for $x_i\in L$, as correctness of the reduction would imply $K(y)\leq m/2-g(m)$ with probability $1-1/n^{\omega(1)}$ (as we could answer NO for all queries $y$ with $K(y)>m/2-g(m)$).
In the second case, the random variable $f(x_i)$ has entropy at most $m/2-g(m)/2+1$. We show that in this case, with probability $1/n^{\Omega(1)}$, $y$ sampled from $f(x_i)$ has $K(y)<m/2$, which again contradicts the correctness of the reduction.
We show that with probability $1/n^{\Omega(1)}$ over $y$ sampled from $f(x_i)$, $K(y\mid x_i)<m/2-\omega(\log(m))$, from which the previous line follows using Proposition 11 and the fact that $K(x_i)=O(\log(n_i))=O(\log(m))$. Indeed let $k$ be the entropy of $f(x_i)$ and let $y_1\ldots y_{2^{k+\log(m)}}$ be the first $2^{k+\log(m)}$ strings in the support of $f(x_i)$ in order of decreasing probability, and let $Y$ be the set of these strings. Conditioned on $x_i$, each such string can be described with a program of size $k+\log(m)+O(\log(m))$, which is $m/2-\omega(\log(m))$, using the facts that $k\leq m/2-g(m)/2+1$ and $g(m)=\Omega(\log(m))$. We claim that with probability at least $1/m$, $y$ sampled from $f(x_i)$ belongs to $Y$. Indeed, if not, we have that with probability greater than $1-1/m$, $y$ sampled from $f(x_i)$ has probability at most $2^{-(k+\log(m))}$, and hence the contribution to the entropy of $f(x_i)$ from such strings is at least $(1-1/m)(k+\log(m))>k$, which is impossible.
Thus we have that with probability $1/m$, $y$ sampled from $f(x_i)$ belongs to $Y$, and that each string in $Y$ has Kolmogorov complexity smaller than $m/2$. But this contradicts the correctness of the reduction, as $x_i$ is a NO instance, and hence with probability $1-1/n^{\omega(1)}$, $y$ sampled from $f(x_i)$ should have $K(y)\geq m/2$ (as otherwise we could answer YES for all queries $y$ with $K(y)<m/2-g(m)$).
\end{proof}

\begin{thebibliography}{99}
\bibitem[1]{published-1} Adi Akavia, Oded Goldreich, Shafi Goldwasser, and Dana Moshkovitz. On basing one-way functions on np-hardness. In Proceedings of the 38th Annual ACM Symposium on Theory of Computing, Seattle, WA, USA, May 21-23, 2006, pages 701–710, 2006.
\bibitem[2]{published-2} Eric Allender. The complexity of complexity. In Computability and Complexity - Essays Dedicated to Rodney G. Downey on the Occasion of His 60th Birthday, pages 79–94, 2017.
\bibitem[3]{published-3} Eric Allender. The new complexity landscape around circuit minimization. In Proceedings of 14th International Conference on Language and Automata Theory and Applications, volume 12038 of Lecture Notes in Computer Science, pages 3–16. Springer, 2020.
\bibitem[4]{published-4} Eric Allender, Harry Buhrman, Luke Friedman, and Bruno Loff. Reductions to the set of random strings: The resource-bounded case. Log. Methods Comput. Sci., 10(3), 2014.
\bibitem[5]{published-5} Eric Allender, Harry Buhrman, and Michal Koucký. What can be efficiently reduced to the kolmogorov-random strings? Ann. Pure Appl. Log., 138(1-3):2–19, 2006.
\bibitem[6]{published-6} Eric Allender, Harry Buhrman, Michal Koucký, Dieter van Melkebeek, and Detlef Ronneburger. Power from random strings. SIAM J. Comput., 35(6):1467–1493, 2006.
\bibitem[7]{published-7} Eric Allender, Harry Buhrman, Michal Koucký, Dieter van Melkebeek, and Detlef Ronneburger. Power from random strings. SIAM J. Comput., 35(6):1467–1493, 2006.
\bibitem[8]{published-8} Eric Allender and Bireswar Das. Zero knowledge and circuit minimization. In Symposium on Mathematical Foundations of Computer Science (MFCS), pages 25–32, 2014. 7 Technically speaking, it might be unclear if Lemma 35 applies here, as the sampled strings do not in general have length m, but the proof of Lemma 35 continues to work as long as the sampled strings have length m$\Omega$(1)
\bibitem[9]{published-9} Eric Allender, George Davie, Luke Friedman, Samuel Hopkins, and Iddo Tzameret. Kolmogorov complexity, circuits, and the strength of formal theories of arithmetic. Chic. J. Theor. Comput. Sci., 2013, 2013.
\bibitem[10]{published-10} Eric Allender, Luke Friedman, and William I. Gasarch. Limits on the computational power of random strings. Inf. Comput., 222:80–92, 2013.
\bibitem[11]{published-11} Eric Allender, Joshua A. Grochow, and Cristopher Moore. Graph isomorphism and circuit size. CoRR, abs/1511.08189, 2015. arXiv:1511.08189.
\bibitem[12]{published-12} Eric Allender, Lisa Hellerstein, Paul McCabe, Toniann Pitassi, and Michael E. Saks. Minimizing disjunctive normal form formulas and AC0 circuits given a truth table. SIAM J. Comput., 38(1):63–84, 2008.
\bibitem[13]{published-13} Eric Allender and Shuichi Hirahara. New insights on the (non-)hardness of circuit minimization and related problems. In International Symposium on Mathematical Foundations of Computer Science (MFCS), pages 54:1–54:14, 2017.
\bibitem[14]{published-14} Eric Allender, Dhiraj Holden, and Valentine Kabanets. The minimum oracle circuit size problem. In International Symposium on Theoretical Aspects of Computer Science (STACS), pages 21–33, 2015.
\bibitem[15]{published-15} Eric Allender, Michal Koucký, Detlef Ronneburger, and Sambuddha Roy. The pervasive reach of resource-bounded Kolmogorov complexity in computational complexity theory. J. Comput. Syst. Sci., 77(1):14–40, 2011.
\bibitem[16]{published-16} Luis Filipe Coelho Antunes and Lance Fortnow. Worst-case running times for average-case algorithms. In Proceedings of the 24th Annual IEEE Conference on Computational Complexity, CCC 2009, Paris, France, 15-18 July 2009, pages 298–303. IEEE Computer Society, 2009.
\bibitem[17]{published-17} Benny Applebaum, Boaz Barak, and David Xiao. On basing lower-bounds for learning on worst-case assumptions. In Proceedings of 49th Annual IEEE Symposium on Foundations of Computer Science, pages 211–220, 2008.
\bibitem[18]{published-18} Sanjeev Arora and Boaz Barak. Computational Complexity: A Modern Approach. Cambridge University Press, 1st edition, 2009.
\bibitem[19]{published-19} Sergei Artemenko, Russell Impagliazzo, Valentine Kabanets, and Ronen Shaltiel. Pseudorandomness when the odds are against you. In Proceedings of the 31st Conference on Computational Complexity, volume 50 of LIPIcs, pages 9:1–9:35. Schloss Dagstuhl - Leibniz- Zentrum für Informatik, 2016.
\bibitem[20]{published-20} Andrej Bogdanov and Luca Trevisan. On worst-case to average-case reductions for NP problems. SIAM J. Comput., 36(4):1119–1159, 2006.
\bibitem[21]{published-21} Harry Buhrman, Lance Fortnow, Michal Koucký, and Bruno Loff. Derandomizing from random strings. In Proceedings of the 25th Annual Conference on Computational Complexity, CCC ’10, 2010.
\bibitem[22]{published-22} Marco L. Carmosino, Russell Impagliazzo, Valentine Kabanets, and Antonina Kolokolova. Agnostic learning from tolerant natural proofs. In Approximation, Randomization, and Combinatorial Optimization, pages 35:1–35:19, 2017.
\bibitem[23]{published-23} Joan Feigenbaum and Lance Fortnow. Random-self-reducibility of complete sets. SIAM J. Comput., 22(5):994–1005, 1993.
\bibitem[24]{published-24} Oded Goldreich and Salil P. Vadhan. On the complexity of computational problems regarding distributions. In Oded Goldreich, editor, Studies in Complexity and Cryptography, volume 6650, pages 390–405. Springer, 2011.
\bibitem[25]{published-25} Shafi Goldwasser and Michael Sipser. Private coins versus public coins in interactive proof systems. In Juris Hartmanis, editor, Proceedings of the 18th Annual ACM Symposium on Theory of Computing, pages 59–68. ACM, 1986.
\bibitem[26]{published-26} Dan Gutfreund and Salil P. Vadhan. Limitations of hardness vs. randomness under uniform reductions. In Proceedings of 12th International Workshop on Randomness and Computation, volume 5171 of Lecture Notes in Computer Science, pages 469–482. Springer, 2008.
\bibitem[27]{published-27} Shuichi Hirahara. Non-black-box worst-case to average-case reductions within NP. In Mikkel Thorup, editor, 59th IEEE Annual Symposium on Foundations of Computer Science, FOCS 2018, Paris, France, October 7-9, 2018, pages 247–258. IEEE Computer Society, 2018.
\bibitem[28]{published-28} Shuichi Hirahara. Unexpected hardness results for kolmogorov complexity under uniform reductions. In Proceedings of 52nd Annual Symposium on Theory of Computing, pages 1038–1051, 2020.
\bibitem[29]{published-29} Shuichi Hirahara, Igor Carboni Oliveira, and Rahul Santhanam. Np-hardness of minimum circuit size problem for OR-AND-MOD circuits. In 33rd Computational Complexity Conference, CCC 2018, pages 5:1–5:31, 2018.
\bibitem[30]{published-30} Shuichi Hirahara and Osamu Watanabe. Limits of minimum circuit size problem as oracle. In Conference on Computational Complexity (CCC), pages 18:1–18:20, 2016.
\bibitem[31]{published-31} Shuichi Hirahara and Osamu Watanabe. On nonadaptive security reductions of hitting set generators. In Proceedings of 24th International Conference on Randomization and Computation, volume 176 of LIPIcs, pages 15:1–15:14. Schloss Dagstuhl - Leibniz-Zentrum für Informatik, 2020.
\bibitem[32]{published-32} John M. Hitchcock. Limitations of efficient reducibility to the kolmogorov random strings. Comput., 1(1):39–43, 2012.
\bibitem[33]{published-33} John M. Hitchcock and Aduri Pavan. On the NP-completeness of the minimum circuit size problem. In Conference on Foundation of Software Technology and Theoretical Computer Science (FSTTCS), pages 236–245, 2015.
\bibitem[34]{published-34} Rahul Ilango. Approaching MCSP from above and below: Hardness for a conditional variant and ac\textasciicircum{}0[p]. In Proceedings of 11th Innovations in Theoretical Computer Science Conference, volume 151 of LIPIcs, pages 34:1–34:26. Schloss Dagstuhl - Leibniz-Zentrum für Informatik, 2020.
\bibitem[35]{published-35} Rahul Ilango. Constant depth formula and partial function versions of MCSP are hard. In Proceedings of 61st IEEE Annual Symposium on Foundations of Computer Science, pages 424–433, 2020.
\bibitem[36]{published-36} Rahul Ilango, Bruno Loff, and Igor Carboni Oliveira. Np-hardness of circuit minimization for multi-output functions. In Proceedings of 35th Computational Complexity Conference, volume 169 of LIPIcs, pages 22:1–22:36. Schloss Dagstuhl - Leibniz-Zentrum für Informatik, 2020.
\bibitem[37]{published-37} Valentine Kabanets and Jin-yi Cai. Circuit minimization problem. In Symposium on Theory of Computing (STOC), pages 73–79, 2000.
\bibitem[38]{published-38} Adam R. Klivans and Dieter van Melkebeek. Graph nonisomorphism has subexponential size proofs unless the polynomial-time hierarchy collapses. SIAM J. Comput., 31(5):1501–1526, 2002.
\bibitem[39]{published-39} Ker-I Ko. On the complexity of learning minimum time-bounded turing machines. SIAM Journal on Computing, 20(5):962–986, 1991.
\bibitem[40]{published-40} Leonid Levin. Universal search problems. Problems of Information Transmission, 9(3):115–116, 1973.
\bibitem[41]{published-41} William J. Masek. Some NP-complete set covering problems. Unpublished Manuscript, 1979.
\bibitem[42]{published-42} Cody Murray and Ryan Williams. On the (non) NP-hardness of computing circuit complexity. In Conference on Computational Complexity (CCC), pages 365–380, 2015.
\bibitem[43]{published-43} Cody Murray and Ryan Williams. On the (non) np-hardness of computing circuit complexity. Theory Comput., 13(1):1–22, 2017.
\bibitem[44]{published-44} Michael Saks and Rahul Santhanam. Circuit lower bounds from np-hardness of MCSP under turing reductions. In Shubhangi Saraf, editor, Proceedings of 35th Computational Complexity Conference, volume 169 of LIPIcs, pages 26:1–26:13. Schloss Dagstuhl - Leibniz-Zentrum für Informatik, 2020.
\bibitem[45]{published-45} Rahul Santhanam. Pseudorandomness and the minimum circuit size problem. In Thomas Vidick, editor, 11th Innovations in Theoretical Computer Science Conference, ITCS 2020, January 12-14, 2020, Seattle, Washington, USA, volume 151 of LIPIcs, pages 68:1–68:26. Schloss Dagstuhl - Leibniz-Zentrum für Informatik, 2020.
\bibitem[46]{published-46} Stathis Zachos. Probabilistic quantifiers and games. J. Comput. Syst. Sci., 36(3):433–451, 1988. doi:10.1016/0022-0000(88)90037-2.
\end{thebibliography}
\end{document}
